% !TeX program = xelatex
% 中文施工主文件：基于 Elsevier CAS 双栏模板，定稿前统一使用中文撰写。
\documentclass[a4paper,fleqn]{cas-dc}

% 使用 TeX Live 自带 Fandol 字体，避免依赖本机中文字体。
\usepackage[UTF8,fontset=fandol]{ctex}
\usepackage[authoryear,longnamesfirst]{natbib}
\usepackage{booktabs}
\usepackage{amsmath,amssymb}
\usepackage{graphicx}
\graphicspath{{figs/}}

\begin{document}
\let\WriteBookmarks\relax

\shorttitle{中文施工稿}
\shortauthors{作者待定}
\title[mode=title]{面向专业视觉语言模型的冻结式视觉适配方法}

\author[1]{作者待定}

\affiliation[1]{organization={单位待定},
            city={城市待定},
            country={中国}}

\begin{abstract}
本文档是基于 Elsevier CAS 双栏模板建立的中文施工稿。论文内容、作者信息、摘要、图表和参考文献将在后续阶段逐步补充；英文翻译仅在中文定稿后进行。
\end{abstract}

\begin{highlights}
\item 使用冻结视觉语言模型主干进行专业领域适配。
\item 通过 MMRL 表征与条件路由实现视觉侧增量知识注入。
\item 在中文阶段完成方法、实验和图表的人工审核。
\end{highlights}

\begin{keywords}
视觉语言模型 \sep 参数高效微调 \sep 专业领域适配 \sep 路由适配器
\end{keywords}

\maketitle

\section{引言}
\label{sec:introduction}

中文正文从这里开始。

\section{方法}
\label{sec:method}

中文方法描述从这里开始。

\section{实验}
\label{sec:experiments}

中文实验内容从这里开始。

\section{结论}
\label{sec:conclusion}

中文结论从这里开始。

% 参考文献在导入真实条目并加入 \cite 后启用，避免空稿阶段触发 BibTeX 报错。
% \bibliographystyle{cas-model2-names}
% \bibliography{references}

\end{document}
